\chapter*{Resumen}
\addcontentsline{toc}{chapter}{Resumen}

% REESCRIBIR A MANO POR ALEX Y JACK
% Borrador, no cerrado. En el reporte del 30 % esta página salió casi entera.
% Léanlo en voz alta y cambien cualquier frase que no dirían en la sustentación.

Queremos saber, \textbf{antes} de integrar, si el YAML de una API bancaria cuenta el mismo negocio
que el \ServiceDomain{} BIAN con el que se la empareja. Hoy esa pregunta se responde a ojo, o no se
responde.

En la mesa abrimos el archivo y el SD. Revisamos los \textit{paths}, el \texttt{operationId} de
cada operación y los \texttt{\$ref} de los \textit{schemas}. Si el contrato dice
\texttt{paymentRequest} y el dominio habla de \texttt{PaymentOrder}, o si un \Behavior{} no tiene
operación en el YAML, no queda un número en el expediente. Sobre esa instancia calculamos similitud
de significado, correspondencia de la jerarquía REST y cobertura de conceptos BIAN:
\SemanticScore{}, \StructuralScore{} y \CoverageScore{}. Las tres entran en un \AlignmentScore{}.
El orden queda en la matriz del Anexo~C.

Esta tesis incluye el prototipo que calcula ese número y un piloto acotado: siete instancias
de Open Finance Brasil contra Service Domains de BIAN v14. El parser y el matching entran.
Quedan fuera el YAML de un banco emisor, la corrida de las 38 APIs contra los 258 dominios,
y la latencia, el \textit{fuzzing} o la seguridad en producción (§\ref{sec:alcances}).

\vspace{0.5cm}
\noindent\textbf{Palabras clave:} \palabrasClaveES.
